%-------------------------------------------------------------
% CV DE BASE - ANGLAIS (modèle vierge)
% Configuration (une seule fois) : l'agent remplace chaque [[MOI: ...]]
% avec le vrai CV de l'utilisateur (traduit en anglais), en gardant
% exactement les mêmes entrées et puces que cv_fr.tex. 1 page.
% Ensuite, ne jamais modifier ce fichier pour une offre : le skill le
% copie dans candidatures/<dossier>/cv.tex et modifie la copie.
% FORMAT FIGÉ : ne rien changer entre \documentclass et
% \begin{document} (polices, marges, espacements, commandes).
%-------------------------------------------------------------
\documentclass[letterpaper,10pt]{article}

\usepackage[letterpaper,left=0.6in,right=0.6in,top=0.42in,bottom=0.4in]{geometry}
\usepackage[T1]{fontenc}
\usepackage[utf8]{inputenc}
\usepackage[english]{babel}
\usepackage{enumitem}
\usepackage[hidelinks]{hyperref}
\usepackage{titlesec}
\input{glyphtounicode}
\pdfgentounicode=1

\pagestyle{empty}
\setlength{\parindent}{0pt}

\titleformat{\section}{\large\scshape}{}{0em}{}[\titlerule]
\titlespacing*{\section}{0pt}{15pt}{6pt}

\setlist[itemize]{leftmargin=1.15em, itemsep=2pt, topsep=5pt, parsep=0pt, partopsep=0pt}

%----------COMMANDES----------
% \experience{Entreprise}{Dates}{Poste}{Lieu}
\newcommand{\experience}[4]{\textbf{#1}\hfill #2\\ \textit{#3}\hfill\textit{#4}}
% \projet{Nom -- Description}{(technos ou contexte)}{Dates}
\newcommand{\projet}[3]{\textbf{#1} {\small #2}\hfill #3}
% espace entre deux entrées
\newcommand{\entre}{\vspace{8pt}}

%=============================================================
\begin{document}

%----------EN-TÊTE----------
\begin{center}
    {\Huge \textbf{[[MOI: First Last]]}} \\[3pt]
    \href{mailto:[[MOI: email]]}{[[MOI: email]]}\enspace|\enspace{}[[MOI: phone]]\enspace|\enspace\href{[[MOI: LinkedIn URL]]}{LinkedIn}\enspace|\enspace\href{[[MOI: GitHub URL]]}{GitHub}
\end{center}
\vspace{-6pt}

%----------EXPERIENCE----------
\section{Experience}

\experience{[[MOI: Company 1]]}{[[MOI: Month YYYY -- Month YYYY]]}{[[MOI: Job title]]}{[[MOI: City, Province]]}
\begin{itemize}
  \item [[MOI: Action verb + what you did + result (with a number if possible)]]
  \item [[MOI: Achievement 2]]
  \item [[MOI: Achievement 3]]
  \item [[MOI: Achievement 4]]
\end{itemize}
\entre

\experience{[[MOI: Company or club 2]]}{[[MOI: Month YYYY -- Present]]}{[[MOI: Role]]}{[[MOI: City, Province]]}
\begin{itemize}
  \item [[MOI: Achievement 1]]
  \item [[MOI: Achievement 2]]
  \item [[MOI: Achievement 3]]
\end{itemize}

%----------PROJECTS----------
\section{Projects}

\projet{[[MOI: Project 1 -- Short description]]}{([[MOI: main technologies]])}{[[MOI: Month YYYY -- Month YYYY]]}
\begin{itemize}
  \item [[MOI: What the project does and your role]]
  \item [[MOI: Technical challenge solved and how]]
  \item [[MOI: Result or impact]]
\end{itemize}
\entre

\projet{[[MOI: Project 2 -- Short description]]}{([[MOI: technologies or context, e.g. hackathon, team of 3]])}{[[MOI: Month YYYY]]}
\begin{itemize}
  \item [[MOI: Achievement 1]]
  \item [[MOI: Achievement 2]]
  \item [[MOI: Achievement 3]]
\end{itemize}
\entre

\projet{[[MOI: Project 3 -- Short description]]}{}{[[MOI: Month YYYY -- Month YYYY]]}
\begin{itemize}
  \item [[MOI: Achievement 1]]
  \item [[MOI: Achievement 2]]
  \item [[MOI: Achievement 3]]
\end{itemize}

%----------EDUCATION----------
\section{Education}

\textbf{[[MOI: School]]}\hfill [[MOI: City, Province]]\\
\textit{[[MOI: Program, e.g. Bachelor's in Software Engineering]]}\hfill\textit{[[MOI: Expected graduation, Month YYYY]]}\\
[[MOI: Credits completed, GPA, awards… or delete this line]]

%----------TECHNICAL SKILLS----------
\section{Technical Skills}

\textbf{Languages:} [[MOI: Python, Java, ...]]\\[2.5pt]
\textbf{Frameworks \& Libraries:} [[MOI: React, Node.js, ...]]\\[2.5pt]
\textbf{Tools:} [[MOI: Git, Linux, Docker, ...]]\\[2.5pt]
\textbf{Methodologies:} [[MOI: Agile/Scrum, unit testing, ...]]

\end{document}
